@section Cursors
The @code{Cursor} class (as well as its subclasses) are used to represent
cursor objects.  This class is never used by an application.  It is used
to house the functionality common to its subclasses.  Therefore, this
class documents functionality which is common to all of its subclasses.

Note that a more convenient mechanism for using cursors is provided by
@code{LoadCursor::Window} and @code{LoadSystemCursor::Window}









@deffn {Copy} Copy::Cursor
@sp 2
@example
@group
c = gCopy(csr);

object  csr;    /*  cursor object          */
object  c;      /*  copy of cursor object  */
@end group
@end example
This method is implemented by all subclasses of @code{Cursor}
(@code{SystemCursor}, @code{ExternalCursor}).

This method is used to create a new cursor object which is a copy of a
given cursor object.  The value of this is to be able to associate a
cursor with multiple objects which will automatically dispose of the
cursor when they are disposed.  If copies weren't used the second object
referencing the cursor would reference a disposed cursor object.

Each cursor object must be disposed (either automatically or manually)
when it is no longer needed.
@example
@group
@exdent Example:

object  c1, c2;

c2 = gCopy(c1);
@end group
@end example
@sp 1
See also:  @code{DeepCopy}
@end deffn










@deffn {DeepCopy} DeepCopy::Cursor
@sp 2
This method performs the same function as @code{Copy}.  See that
method for details.
@end deffn










@deffn {DeepDispose} DeepDispose::Cursor
@sp 2
This method performs the same function as @code{Dispose}.  See that
method for details.
@end deffn







@deffn {Dispose} Dispose::Cursor
@sp 2
@example
@group
r = gDispose(csr);

object  csr;   /*  a cursor object    */
object  r;     /*  NULL               */
@end group
@end example
This method is used to dispose of a cursor object when it
is no longer needed.  This method is rarely needed due to the fact that
when a window is disposed it automatically calls this method to
dispose of its associated cursor object.

The value returned is always @code{NULL} and may be used to null out
the variable which contained the object being disposed in order to
avoid future accidental use.
@example
@group
@exdent Example:

object  csr;

csr = gDispose(csr);
@end group
@end example
@sp 1
See also:  @code{DeepDispose}
@end deffn










@deffn {Handle} Handle::Cursor
@sp 2
@example
@group
h = gHandle(csr);

object  csr;    /*  cursor object   */
HANDLE  h;      /*  Windows handle  */
@end group
@end example
This method is used to obtain the Windows internal handle associated with
a cursor object.  

Note that this method may be used with most WDS objects in order to obtain
the internal handle that Windows normally associates with each type of object.
See the appropriate documentation.
@example
@group
@exdent Example:

object  csr;
HCURSOR h;

h = gHandle(csr);
@end group
@end example
@c @sp 1
@c See also:  @code{}
@end deffn









@subsection System Cursors
The @code{SystemCursor} class represents Windows predefined cursors.
This class is a subclass of @code{Cursor} and as such inherits most
of its functionality from it.  This section documents the methods
particular to this class.  See the @code{Cursor} class for additional
functionality.











@deffn {LoadSys} LoadSys::SystemCursor
@sp 2
@example
@group
csr = gLoadSys(SystemCursor, id);

char    *id;    /*  cursor id      */
object  csr;    /*  cursor object  */
@end group
@end example
This method is used to create a new object which represents one of the
Windows predefined cursors.  @code{id} should be one of the Windows
defined macros specifying the desired cursor.  It is defined in the
Windows documentation under the function @code{LoadCursor} and
starts with @code{IDC_}.
@example
@group
@exdent Example:

object  csr;

csr = gLoadSys(SystemCursor, IDC_IBEAM);
@end group
@end example
@sp 1
See also:  @code{LoadSystemCursor::Window, SetCursor::Application}
@end deffn





@subsection External Cursors
The @code{ExternalCursor} class represents arbitrary application defined
cursors.  This class is a subclass of @code{Cursor} and as such inherits
most of its functionality from it.  This section documents the methods
particular to this class.  See the @code{Cursor} class for additional
functionality.






@deffn {Load} Load::ExternalCursor
@sp 2
@example
@group
csr = mLoad(ExternalCursor, id);

unsigned id;    /*  cursor id      */
object   csr;   /*  cursor object  */
@end group
@end example
This method is used to load a programmer defined application specific
cursor.

@code{id} is a programmer defined unsigned integer which identifies
the cursor.  This identifier is normally a macro and defined through the
resource editor.  

The value returned is an object representing the cursor loaded, or
@code{NULL} if the cursor was not found.
@example
@group
@exdent Example:

object  csr;

csr = mLoad(ExternalCursor, MY_CURSOR);
@end group
@end example
@sp 1
See also:  @code{LoadCursor::Window, SetCursor::Application}
@end deffn











@deffn {LoadStr} LoadStr::ExternalCursor
@sp 2
@example
@group
csr = mLoadStr(ExternalCursor, id);

char    *id;    /*  cursor id      */
object   csr;   /*  cursor object  */
@end group
@end example
This method is used to load a programmer defined application specific
cursor by name.

@code{id} is a programmer defined name which identifies the cursor.
This identifier is normally defined through the resource editor.

The value returned is an object representing the cursor loaded, or
@code{NULL} if the cursor was not found.
@example
@group
@exdent Example:

object  csr;

csr = mLoadStr(ExternalCursor, "mycursor");
@end group
@end example
@sp 1
See also:  @code{LoadCursor::Window, SetCursor::Application,}
        @code{Use::Window}
@end deffn




@section Icons
The @code{Icon} class (as well as its subclasses) are used to represent
icon objects.  This class is never used by an application.  It is used
to house the functionality common to its subclasses.  Therefore, this
class documents functionality which is common to all of its subclasses.

Note that a more convenient mechanism for using icons is provided by
@code{LoadIcon::Window} and @code{LoadSystemIcon::Window}









@deffn {Copy} Copy::Icon
@sp 2
@example
@group
c = gCopy(icn);

object  icn;    /*  icon object          */
object  c;      /*  copy of icon object  */
@end group
@end example
This method is implemented by all subclasses of @code{Icon}
(@code{SystemIcon}, @code{ExternalIcon}).

This method is used to create a new icon object which is a copy of a
given icon object.  The value of this is to be able to associate a
icon with multiple objects which will automatically dispose of the
icon when they are disposed.  If copies weren't used the second object
referencing the icon would reference a disposed icon object.

Each icon object must be disposed (either automatically or manually)
when it is no longer needed.
@example
@group
@exdent Example:

object  i1, i2;

i2 = gCopy(i1);
@end group
@end example
@sp 1
See also:  @code{DeepCopy}
@end deffn










@deffn {DeepCopy} DeepCopy::Icon
@sp 2
This method performs the same function as @code{Copy}.  See that
method for details.
@end deffn










@deffn {DeepDispose} DeepDispose::Icon
@sp 2
This method performs the same function as @code{Dispose}.  See that
method for details.
@end deffn







@deffn {Dispose} Dispose::Icon
@sp 2
@example
@group
r = gDispose(icn);

object  icn;   /*  an icon object  */
object  r;     /*  NULL            */
@end group
@end example
This method is used to dispose of an icon object when it is no longer
needed.  This method is rarely needed due to the fact that when a window
is disposed it automatically calls this method to dispose of its
associated icon object.

The value returned is always @code{NULL} and may be used to null out
the variable which contained the object being disposed in order to
avoid future accidental use.
@example
@group
@exdent Example:

object  icn;

icn = gDispose(icn);
@end group
@end example
@sp 1
See also:  @code{DeepDispose}
@end deffn










@deffn {Handle} Handle::Icon
@sp 2
@example
@group
h = gHandle(icn);

object  icn;    /*  icon object     */
HANDLE  h;      /*  Windows handle  */
@end group
@end example
This method is used to obtain the Windows internal handle associated with
an icon object.  

Note that this method may be used with most WDS objects in order to obtain
the internal handle that Windows normally associates with each type of object.
See the appropriate documentation.
@example
@group
@exdent Example:

object  icn;
HICON   h;

h = gHandle(icn);
@end group
@end example
@c @sp 1
@c See also:  @code{}
@end deffn









@subsection System Icons
The @code{SystemIcon} class represents Windows predefined icons.
This class is a subclass of @code{Icon} and as such inherits most
of its functionality from it.  This section documents the methods
particular to this class.  See the @code{Icon} class for additional
functionality.











@deffn {LoadSys} LoadSys::SystemIcon
@sp 2
@example
@group
icn = gLoadSys(SystemIcon, id);

char    *id;    /*  icon id      */
object  icn;    /*  icon object  */
@end group
@end example
This method is used to create a new object which represents one of the
Windows predefined icons.  @code{id} should be one of the Windows
defined macros specifying the desired icon.  It is defined in the
Windows documentation under the function @code{LoadIcon} and
starts with @code{IDI_}.
@example
@group
@exdent Example:

object  icn;

icn = gLoadSys(SystemIcon, IDI_HAND);
@end group
@end example
@sp 1
See also:  @code{LoadSystemIcon::Window, SetIcon::Application}
@end deffn





@subsection External Icons
The @code{ExternalIcon} class represents arbitrary application defined
icons.  This class is a subclass of @code{Icon} and as such inherits
most of its functionality from it.  This section documents the methods
particular to this class.  See the @code{Icon} class for additional
functionality.






@deffn {Load} Load::ExternalIcon
@sp 2
@example
@group
icn = mLoad(ExternalIcon, id);

unsigned id;    /*  icon id      */
object   icn;   /*  icon object  */
@end group
@end example
This method is used to load a programmer defined application specific
icon.

@code{id} is a programmer defined unsigned integer which identifies
the icon.  This identifier is normally a macro and defined through the
resource editor.  

The value returned is an object representing the icon loaded, or
@code{NULL} if the icon was not found.
@example
@group
@exdent Example:

object  icn;

icn = mLoad(ExternalIcon, MY_ICON);
@end group
@end example
@sp 1
See also:  @code{LoadIcon::Window, SetIcon::Application}
@end deffn











@deffn {LoadStr} LoadStr::ExternalIcon
@sp 2
@example
@group
icn = mLoadStr(ExternalIcon, id);

char    *id;    /*  icon id      */
object   icn;   /*  icon object  */
@end group
@end example
This method is used to load a programmer defined application specific
icon by name.

@code{id} is a programmer defined name which identifies the icon.
This identifier is normally defined through the resource editor.

The value returned is an object representing the icon loaded, or
@code{NULL} if the icon was not found.
@example
@group
@exdent Example:

object  icn;

icn = mLoadStr(ExternalIcon, "myicon");
@end group
@end example
@sp 1
See also:  @code{LoadIcon::Window, SetIcon::Application,}
        @code{Use::Window}
@end deffn





@section Fonts
The @code{Font} class (as well as its subclasses) are used to represent
font objects.  This class is never used by an application.  It is used
to house the functionality common to its subclasses.  Therefore, this
class documents functionality which is common to all of its subclasses.

Note that a more convenient mechanism for using fonts is provided by
@code{LoadFont::Window, LoadSystemFont::Window, LoadFont::Printer},
and @code{LoadSystemFont::Printer}.










@deffn {AveCharWidth} AveCharWidth::Font
@sp 2
@example
@group
r = gAveCharWidth(fnt);

object   fnt;  /*  a font object  */
int      r;    /*  ave char width */
@end group
@end example
This method is used to gain access to the average character width
associated with a particular font.  
@example
@group
@exdent Example:

object  fnt;
int     cw;

cw = gAveCharWidth(fnt);
@end group
@end example
@sp 1
See also:  @code{LineHeight, GetTM}
@end deffn










@deffn {Copy} Copy::Font
@sp 2
@example
@group
c = gCopy(fnt);

object  fnt;    /*  font object          */
object  c;      /*  copy of font object  */
@end group
@end example
This method is implemented by all subclasses of @code{Font}
(@code{SystemFont}, @code{ExternalFont}).

This method is used to create a new font object which is a copy of a
given font object.  The value of this is to be able to associate a
font with multiple objects which will automatically dispose of the
font when they are disposed.  If copies weren't used the second object
referencing the font would reference a disposed font object.

Each font object must be disposed (either automatically or manually)
when it is no longer needed.
@example
@group
@exdent Example:

object  f1, f2;

f2 = gCopy(f1);
@end group
@end example
@sp 1
See also:  @code{DeepCopy}
@end deffn











@deffn {DeepCopy} DeepCopy::Font
@sp 2
This method performs the same function as @code{Copy}.  See that
method for details.
@end deffn











@deffn {DeepDispose} DeepDispose::Font
@sp 2
This method performs the same function as @code{Dispose}.  See that
method for details.
@end deffn








@deffn {Dispose} Dispose::Font
@sp 2
@example
@group
r = gDispose(fnt);

object  fnt;   /*  a font object  */
object  r;     /*  NULL           */
@end group
@end example
This method is used to dispose of a font object when it is no longer
needed.  This method is rarely needed due to the fact that when a window
or printer object is disposed it automatically calls this method to
dispose of its associated font object.

The value returned is always @code{NULL} and may be used to null out
the variable which contained the object being disposed in order to
avoid future accidental use.
@example
@group
@exdent Example:

object  fnt;

fnt = gDispose(fnt);
@end group
@end example
@sp 1
See also:  @code{DeepDispose}
@end deffn










@deffn {GetTM} GetTM::Font
@sp 2
@example
@group
r = gGetTM(fnt, tm);

object      fnt;  /*  a font object  */
TEXTMETRIC  *tm;  /*  font metrics   */
object      r;    /*  fnt            */
@end group
@end example
This method is used to gain access to the font metrics associated with
a particular font.
@example
@group
@exdent Example:

object  fnt;
TEXTMETRIC  tm;

gGetTM(fnt, &tm);
@end group
@end example
@sp 1
See also:  @code{LineHeight, AveCharWidth}
@end deffn













@deffn {Handle} Handle::Font
@sp 2
@example
@group
h = gHandle(fnt);

object  fnt;    /*  font object     */
HANDLE  h;      /*  Windows handle  */
@end group
@end example
This method is used to obtain the Windows internal handle associated with
a font object.  

Note that this method may be used with most WDS objects in order to obtain
the internal handle that Windows normally associates with each type of object.
See the appropriate documentation.
@example
@group
@exdent Example:

object  fnt;
HGDIOBJ h;

h = gHandle(fnt);
@end group
@end example
@c @sp 1
@c See also:  @code{}
@end deffn














@deffn {LineHeight} LineHeight::Font
@sp 2
@example
@group
r = gLineHeight(fnt);

object   fnt;  /*  a font object  */
int      r;    /*  line height    */
@end group
@end example
This method is used to gain access to the line height associated with
a particular font.  This number is equal to the height of the largest
character in the font plus a reasonable amount of space to separate
lines containing the font.
@example
@group
@exdent Example:

object  fnt;
int     lh;

lh = gLineHeight(fnt);
@end group
@end example
@sp 1
See also:  @code{AveCharWidth, GetTM}
@end deffn




@deffn {Name} Name::Font
@sp 2
@example
@group
s = gName(fnt);

object  fnt;  /*  a font object  */
char    *s;   /*  font face name */
@end group
@end example
This instance method returns the font's face name including style
suffixes (e.g., "Arial Bold Italic").
@example
@group
@exdent Example:

object  fnt;
char    *name;

name = gName(fnt);
@end group
@end example
@end deffn




@deffn {PixelHeight} PixelHeight::Font
@sp 2
@example
@group
n = gPixelHeight(Font);

int  n;  /*  vertical DPI (pixels per logical inch)  */
@end group
@end example
This class method returns the screen's vertical DPI (pixels per
logical inch).  It is called on the @code{Font} class itself, not on
an instance.
@example
@group
@exdent Example:

int  dpi;

dpi = gPixelHeight(Font);
@end group
@end example
@end deffn




@deffn {PointSize} PointSize::Font
@sp 2
@example
@group
n = gPointSize(fnt);

object  fnt;  /*  a font object  */
int     n;    /*  point size     */
@end group
@end example
This instance method returns the font's point size.
@example
@group
@exdent Example:

object  fnt;
int     ps;

ps = gPointSize(fnt);
@end group
@end example
@end deffn








@subsection System Fonts
The @code{SystemFont} class represents Windows predefined fonts.
This class is a subclass of @code{Font} and as such inherits most
of its functionality from it.  This section documents the methods
particular to this class.  See the @code{Font} class for additional
functionality.








@deffn {Load} Load::SystemFont
@sp 2
@example
@group
fnt = mLoad(SystemFont, id);

unsigned id;    /*  icon id      */
object   fnt;   /*  icon object  */
@end group
@end example
This method is used to create a new object which represents one of the
Windows predefined fonts.  @code{id} should be one of the Windows
defined macros specifying the desired font.  It is defined in the
Windows documentation under the function @code{GetStockObject} and
ends with @code{_FONT}.
@example
@group
@exdent Example:

object  fnt;

fnt = mLoad(SystemFont, ANSI_FIXED_FONT);
@end group
@end example
@sp 1
See also:  @code{LoadSystemFont::Window, LoadSystemFont::Printer,}
@iftex
@hfil @break @hglue .63in 
@end iftex
@code{SetFont::Application, Use::Window, Use::Printer}
@end deffn











@subsection External Fonts
The @code{ExternalFont} class represents arbitrary external fonts.  This
class is a subclass of @code{Font} and as such inherits most of its
functionality from it.  This section documents the methods particular to
this class.  See the @code{Font} class for additional functionality.









@deffn {Indirect} Indirect::ExternalFont
@sp 2
@example
@group
fnt = gIndirect(ExternalFont, lf);

LOGFONT *lf;    /*  logical font  */
object  fnt;    /*  font object   */
@end group
@end example
This method is used to find and load a font which most closely matches
the parameters indicated by @code{lf}.  The structure of @code{lf} is
fully described in the Windows documentation under the structure
@code{LOGFONT}.

The value returned is an object representing the font loaded, or
@code{NULL} if the font was not found.
@example
@group
@exdent Example:

LOGFONT lf;
object  fnt;

/*  initialize lf  */
fnt = gIndirect(ExternalFont, &lf);
@end group
@end example
@sp 1
See also:  @code{LoadFont::Window, LoadFont::Printer, SetFont::Application,}
@iftex
@hfil @break @hglue .63in 
@end iftex
@code{Use::Window, Use::Printer}
@end deffn










@deffn {New} New::ExternalFont
@sp 2
@example
@group
fnt = vNew(ExternalFont, nm, ps);

char    *nm;    /*  font name    */
int     ps;     /*  point size   */
object  fnt;    /*  font object  */
@end group
@end example
This method is used to load an arbitrary font by name at any point size.
@code{nm} is the full name of the font as it appears
when you list the available fonts via the control-panel / fonts Windows
utility, minus the font type in parentheses.  @code{ps} indicates the
desired point size.

The value returned is an object representing the font loaded, or
@code{NULL} if the font was not found.

Note that @code{nm} may also be a Dynace object cast as a (@code{char *}).
@example
@group
@exdent Example:

object  fnt;

fnt = vNew(ExternalFont, "Times New Roman", 12);
@end group
@end example
@sp 1
See also:  @code{LoadFont::Window, LoadFont::Printer, SetFont::Application,}
@iftex
@hfil @break @hglue .63in 
@end iftex
@code{Use::Window, Use::Printer}
@end deffn









@section Brushes
The @code{Brush} class (as well as its subclasses) are used to represent
brush objects.  This class is never used by an application.  It is used
to house the functionality common to its subclasses.  Therefore, this
class documents functionality which is common to all of its subclasses.









@deffn {Color} Color::Brush
@sp 2
@example
@group
c = gColor(bsh);

object    bsh;  /*  brush object   */
COLORREF  c;    /*  brush color    */
@end group
@end example
This method is used to obtain the color associated with a brush.
@example
@group
@exdent Example:

object  bsh;
COLORREF  c;

c = gColor(bsh);
@end group
@end example
@c @sp 1
@c See also:  @code{}
@end deffn










@deffn {Copy} Copy::Brush
@sp 2
@example
@group
c = gCopy(bsh);

object  bsh;    /*  brush object          */
object  c;      /*  copy of brush object  */
@end group
@end example
This method is implemented by all subclasses of @code{Brush}
(@code{SystemBrush}, @code{StockBrush}, @code{SolidBrush},
@code{HatchBrush}).

This method is used to create a new brush object which is a copy of a
given brush object.  The value of this is to be able to associate a
brush with multiple objects which will automatically dispose of the
brush when they are disposed.  If copies weren't used the second object
referencing the brush would reference a disposed brush object.

Each brush object must be disposed (either automatically or manually)
when it is no longer needed.
@example
@group
@exdent Example:

object  b1, b2;

b2 = gCopy(b1);
@end group
@end example
@sp 1
See also:  @code{DeepCopy}
@end deffn









@deffn {DeepCopy} DeepCopy::Brush
@sp 2
This method performs the same function as @code{Copy}.  See that
method for details.
@end deffn





@deffn {DeepDispose} DeepDispose::Brush
@sp 2
This method performs the same function as @code{Dispose}.  See that
method for details.
@end deffn







@deffn {Dispose} Dispose::Brush
@sp 2
@example
@group
r = gDispose(bsh);

object  bsh;   /*  a brush object  */
object  r;     /*  NULL            */
@end group
@end example
This method is used to dispose of a brush object when it is no longer
needed.  This method is rarely needed due to the fact that when a window
or printer object is disposed it automatically calls this method to
dispose of its associated brush object.

The value returned is always @code{NULL} and may be used to null out
the variable which contained the object being disposed in order to
avoid future accidental use.
@example
@group
@exdent Example:

object  bsh;

bsh = gDispose(bsh);
@end group
@end example
@sp 1
See also:  @code{DeepDispose}
@end deffn






@deffn {Handle} Handle::Brush
@sp 2
@example
@group
h = gHandle(bsh);

object  bsh;    /*  brush object    */
HANDLE  h;      /*  Windows handle  */
@end group
@end example
This method is used to obtain the Windows internal handle associated with
a brush object.  

Note that this method may be used with most WDS objects in order to obtain
the internal handle that Windows normally associates with each type of object.
See the appropriate documentation.
@example
@group
@exdent Example:

object  bsh;
HBRUSH  h;

h = gHandle(bsh);
@end group
@end example
@c @sp 1
@c See also:  @code{}
@end deffn








@subsection Stock Brushes
The @code{StockBrush} class represents common Windows defined brushes.  This
class is a subclass of @code{Brush} and as such inherits most of its
functionality from it.  This section documents the methods particular to
this class.  See the @code{Brush} class for additional functionality.






@deffn {New} New::StockBrush
@sp 2
@example
@group
bsh = vNew(StockBrush, id);

unsigned id;    /*  brush id      */
object  bsh;    /*  brush object  */
@end group
@end example
This method is used to load one of the Windows pre-defined brushes.
The available options are macros documented in the Windows documentation
under the function @code{GetStockObject} and end with @code{_BRUSH}.

The value returned is an object representing the brush loaded, or
@code{NULL} if the brush was not found.
@example
@group
@exdent Example:

object  bsh;

bsh = vNew(StockBrush, GRAY_BRUSH);
@end group
@end example
@sp 1
See also:  @code{TextBrush::Window, BackBrush::Window,}
        @code{Use::Window, Use::Printer}
@end deffn









@subsection Solid Brushes
The @code{SolidBrush} class represents arbitrary application defined
brushes.  This class is a subclass of @code{Brush} and as such inherits
most of its functionality from it.  This section documents the methods
particular to this class.  See the @code{Brush} class for additional
functionality.






@deffn {New} New::SolidBrush
@sp 2
@example
@group
bsh = vNew(SolidBrush, red, green, blue);

int     red;
int     green;
int     blue;
object  bsh;    /*  brush object  */
@end group
@end example
This method is used to create a solid colored brush with specific colors.
Each color argument specifies a number between 0 and 255 and indicates
the intensity of the related color.  If all three are 0 you get black,
and all three at 255 is white.

The value returned is an object representing the brush created, or
@code{NULL} if the brush was not created.
@example
@group
@exdent Example:

object  bsh;

bsh = vNew(SolidBrush, 255, 0, 0);
@end group
@end example
@sp 1
See also:  @code{TextBrush::Window, BackBrush::Window,}
        @code{Use::Window, Use::Printer}
@end deffn




@subsection System Brushes
The @code{SystemBrush} class represents brushes which are those colors
which were selected by the user as global to their Windows environment.
The programmer may select, for example, the color the user chose for
background windows.

This class is a subclass of @code{Brush} and as such inherits
most of its functionality from it.  This section documents the methods
particular to this class.  See the @code{Brush} class for additional
functionality.








@deffn {New} New::SystemBrush
@sp 2
@example
@group
bsh = vNew(SystemBrush, id);

unsigned id;    /*  brush id      */
object  bsh;    /*  brush object  */
@end group
@end example
This method is used to load one of the Windows user defined brushes.
The available options are macros documented in the Windows documentation
under the function @code{GetSysColor}.

The value returned is an object representing the brush loaded, or
@code{NULL} if the brush was not found.
@example
@group
@exdent Example:

object  bsh;

bsh = vNew(SystemBrush, COLOR_ACTIVEBORDER);
@end group
@end example
@sp 1
See also:  @code{TextBrush::Window, BackBrush::Window,}
        @code{Use::Window, Use::Printer}
@end deffn










@subsection Hatch Brushes
The @code{HatchBrush} class represents arbitrary application defined
brushes with a specified pattern.  This class is a subclass of
@code{Brush} and as such inherits most of its functionality from it.
This section documents the methods particular to this class.  See the
@code{Brush} class for additional functionality.






@deffn {New} New::HatchBrush
@sp 2
@example
@group
bsh = vNew(HatchBrush, red, green, blue, style);

int     red;
int     green;
int     blue;
int     style;  /*  brush pattern  */
object  bsh;    /*  brush object   */
@end group
@end example
This method is used to create a colored brush with specific colors and
pattern Each color argument specifies a number between 0 and 255 and
indicates the intensity of the related color.  If all three are 0 you
get black, and all three at 255 is white.

The @code{style} parameter indicates the specific pattern for the brush.
This parameter is a macro documented in the Windows documentation under
the function @code{CreateHatchBrush} and begin with @code{HS_}.

The value returned is an object representing the brush created, or
@code{NULL} if the brush was not created.
@example
@group
@exdent Example:

object  bsh;

bsh = vNew(HatchBrush, 255, 10, 10, HS_VERTICAL);
@end group
@end example
@sp 1
See also:  @code{TextBrush::Window, BackBrush::Window,}
        @code{Use::Window, Use::Printer}
@end deffn









@section Pens
The @code{Pen} class (as well as its subclasses) are used to represent
pen objects.  This class is never used by an application.  It is used
to house the functionality common to its subclasses.  Therefore, this
class documents functionality which is common to all of its subclasses.









@deffn {Color} Color::Pen
@sp 2
@example
@group
c = gColor(pn);

object    pn;   /*  pen object   */
COLORREF  c;    /*  pen color    */
@end group
@end example
This method is used to obtain the color associated with a pen.
@example
@group
@exdent Example:

object    pn;
COLORREF  c;

c = gColor(pn);
@end group
@end example
@c @sp 1
@c See also:  @code{}
@end deffn










@deffn {Copy} Copy::Pen
@sp 2
@example
@group
c = gCopy(pn);

object  pn;     /*  pen object          */
object  c;      /*  copy of pen object  */
@end group
@end example
This method is implemented by all subclasses of @code{Pen}
(@code{StockPen}, @code{CustomPen}).

This method is used to create a new pen object which is a copy of a
given pen object.  The value of this is to be able to associate a
pen with multiple objects which will automatically dispose of the
pen when they are disposed.  If copies weren't used the second object
referencing the pen would reference a disposed pen object.

Each pen object must be disposed (either automatically or manually)
when it is no longer needed.
@example
@group
@exdent Example:

object  p1, p2;

p2 = gCopy(p1);
@end group
@end example
@sp 1
See also:  @code{DeepCopy}
@end deffn









@deffn {DeepCopy} DeepCopy::Pen
@sp 2
This method performs the same function as @code{Copy}.  See that
method for details.
@end deffn





@deffn {DeepDispose} DeepDispose::Pen
@sp 2
This method performs the same function as @code{Dispose}.  See that
method for details.
@end deffn







@deffn {Dispose} Dispose::Pen
@sp 2
@example
@group
r = gDispose(pn);

object  pn;   /*  a pen object  */
object  r;    /*  NULL          */
@end group
@end example
This method is used to dispose of a pen object when it is no longer
needed.  This method is rarely needed due to the fact that when a window
or printer object is disposed it automatically calls this method to
dispose of its associated pen object.

The value returned is always @code{NULL} and may be used to null out
the variable which contained the object being disposed in order to
avoid future accidental use.
@example
@group
@exdent Example:

object  pn;

pn = gDispose(pn);
@end group
@end example
@sp 1
See also:  @code{DeepDispose}
@end deffn






@deffn {Handle} Handle::Pen
@sp 2
@example
@group
h = gHandle(pn);

object  pn;     /*  pen object    */
HANDLE  h;      /*  Windows handle  */
@end group
@end example
This method is used to obtain the Windows internal handle associated with
a pen object.  

Note that this method may be used with most WDS objects in order to obtain
the internal handle that Windows normally associates with each type of object.
See the appropriate documentation.
@example
@group
@exdent Example:

object  pn;
HPEN    h;

h = gHandle(pn);
@end group
@end example
@c @sp 1
@c See also:  @code{}
@end deffn









@subsection Stock Pens
The @code{StockPen} class represents common Windows defined pens.  This
class is a subclass of @code{Pen} and as such inherits most of its
functionality from it.  This section documents the methods particular to
this class.  See the @code{Pen} class for additional functionality.






@deffn {New} New::StockPen
@sp 2
@example
@group
pn = vNew(StockPen, id);

unsigned id;    /*  pen id      */
object  pn;     /*  pen object  */
@end group
@end example
This method is used to load one of the Windows pre-defined pens.
The available options are macros documented in the Windows documentation
under the function @code{GetStockObject} and end with @code{_PEN}.

The value returned is an object representing the pen loaded, or
@code{NULL} if the pen was not found.
@example
@group
@exdent Example:

object  pn;

pn = vNew(StockPen, BLACK_PEN);
@end group
@end example
@sp 1
See also:  @code{Use::Printer}
@end deffn







@subsection Custom Pens
The @code{CustomPen} class represents arbitrary application defined pens
with a specified pattern and width.  This class is a subclass of
@code{Pen} and as such inherits most of its functionality from it.  This
section documents the methods particular to this class.  See the
@code{Pen} class for additional functionality.






@deffn {New} New::CustomPen
@sp 2
@example
@group
pn = vNew(CustomPen, red, green, blue, style, width);

int     red;
int     green;
int     blue;
int     style;  /*  pen pattern     */
int     width;  /*  pen line width  */
object  pn;     /*  pen object      */
@end group
@end example
This method is used to create a colored pen with specific colors,
pattern and line width.  Each color argument specifies a number between
0 and 255 and indicates the intensity of the related color.  If all
three are 0 you get black, and all three at 255 is white.

The @code{style} parameter indicates the specific pattern for the pen.
This parameter is a macro documented in the Windows documentation under
the function @code{CreatePen} and begin with @code{PS_}.
The @code{width} parameter specifies the width of the line in logical
units.

The value returned is an object representing the pen created, or
@code{NULL} if the pen was not created.
@example
@group
@exdent Example:

object  pn;

pn = vNew(CustomPen, 255, 10, 10, PS_DASH, 5);
@end group
@end example
@sp 1
See also:  @code{Use::Printer}
@end deffn




@section Help System
The @code{HelpSystem} class is used to support the standard Windows
help system.  In addition to being able to support arbitrary evocation
of help messages, this class is used by other WDS classes in order to
provide complete support for context sensitive help to any level.
Each class provides mechanisms in order to associate context sensitive
help text to windows, dialogs, menus or controls.

Most classes which provide context sensitive help support have a method
called @code{SetTopic} which is used to associate a particular topic
within a given context.

All the methods in this class are class methods.  This means that
the first argument to all the methods will be @code{HelpSystem} and
there is never an instance object to keep track of.

See the appropriate examples for detailed information on how to
create the actual help file.











@deffn {HelpContents} HelpContents::HelpSystem
@sp 2
@example
@group
r = gHelpContents(HelpSystem);

object  r;      /*  HelpSystem  */
@end group
@end example
This method is used to display the contents screen of the help file.
@code{HelpSystem} will be returned if the request is successful and
@code{NULL} otherwise.
@example
@group
@exdent Example:

gHelpContents(HelpSystem);
@end group
@end example
@sp 1
See also:  @code{HelpFile, HelpTopic, HelpInContext}
@end deffn















@deffn {HelpFile} HelpFile::HelpSystem
@sp 2
@example
@group
r = gHelpFile(HelpSystem, hf);

char    *hf;    /*  help file   */
object  r;      /*  HelpSystem  */
@end group
@end example
This method is used to determine the external file to be used for all help
references.  This must be done prior to any of the other help facilities
use.  The help file will not be opened until one of the help display
methods is evoked.  The application must also have a main window prior
to the evocation of the help system.  The help system will automatically
terminate when the user terminates the application.
@example
@group
@exdent Example:

gHelpFile(HelpSystem, "helpfile.hlp");
@end group
@end example
@sp 1
See also:  @code{HelpContents}
@end deffn














@deffn {HelpInContext} HelpInContext::HelpSystem
@sp 2
@example
@group
r = gHelpInContext(HelpSystem);

object  r;      /*  HelpSystem  */
@end group
@end example
This method is used to display the help screen associated with the
current context (set with @code{SetTopic}).  If no context is set the
help contents will be displayed.  @code{HelpSystem} will be returned if
the request is successful and @code{NULL} otherwise.

Note that @code{SetTopic} and @code{HelpInContext}, although available
to an application program, are mainly used internally by WDS in order to
support the context sensitive help facility provided by the @code{Window},
@code{Dialog}, @code{Menu}, and @code{Control} classes.  See the
@code{SetTopic} method associated with those classes.
@example
@group
@exdent Example:

gHelpInContext(HelpSystem);
@end group
@end example
@sp 1
See also:  @code{HelpFile, SetTopic, HelpContents, HelpTopic}
@end deffn







@deffn {HelpTopic} HelpTopic::HelpSystem
@sp 2
@example
@group
r = gHelpTopic(HelpSystem, tpc);

char    *tpc;   /*  help topic  */
object  r;      /*  HelpSystem  */
@end group
@end example
This method is used to display the help screen associated with a given topic.
@iftex
@hfil @break 
@end iftex
@code{HelpSystem} will be returned if the request is successful
and @code{NULL} otherwise.
@example
@group
@exdent Example:

gHelpTopic(HelpSystem, "fileOpen");
@end group
@end example
@sp 1
See also:  @code{HelpFile, HelpContents, HelpInContext}
@end deffn














@deffn {SetTopic} SetTopic::HelpSystem
@sp 2
@example
@group
r = gSetTopic(HelpSystem, tpc);

char    *tpc;   /*  help topic      */
char    *r;     /*  previous topic  */
@end group
@end example
This method is used to set the current help topic context.  It is
used by @code{HelpInContext} to display the help text associated
with the current context.  The previous help topic will be returned.

Note that @code{SetTopic} and @code{HelpInContext}, although available
to an application program, are mainly used internally by WDS in order to
support the context sensitive help facility provided by the @code{Window},
@code{Dialog}, @code{Menu}, and @code{Control} classes.  See the
@code{SetTopic} method associated with those classes.
@example
@group
@exdent Example:

gSetTopic(HelpSystem, "someTopic");
@end group
@end example
@sp 1
See also:  @code{HelpFile, HelpInContext, HelpContents, HelpTopic}
@end deffn










@section Common Dialogs
The @code{CommonDialog} class is only used to group the common dialogs
which are all subclasses of this class.  There is no specific
functionality associated with this class.  See the particular subclasses
for documentation.






@subsection File Selection Dialog
The @code{FileDialog} class provides a standard and convenient method of
querying the user for a file name.  The user is able to browse the disk
or enter a new file name.

The @code{fd} parameter used by all methods in this class refer to
the file dialog object returned by the @code{New} method.














@deffn {AppendFilter} AppendFilter::FileDialog
@sp 2
@example
@group
r = gAppendFilter(fd, ttl, fltr);

object  fd;     /*  file dialog  */
char    *ttl;   /*  title        */
char    *fltr;  /*  file filter  */
object  r;      /*  file dialog  */
@end group
@end example
This method is used to group files according to some file naming
criteria for user selection.  There may be any number of groups.  The
user selects the group they are interested in and are presented with
a list of files which meet the criteria.

@code{ttl} is the name of the group which the user is presented with.
@code{fltr} is the file filter and may consist of several filters,
each separated with a semicolon.

This method may be called any number of times to create several groups.
@example
@group
@exdent Example:

object  fd;

gAppendFilter(fd, "Document Files", "*.txt;*.doc");
gAppendFilter(fd, "Source Files", "*.c;*.h");
@end group
@end example
@sp 1
See also:  @code{SetFile, InitialDir, DefExt}
@end deffn














@deffn {DefExt} DefExt::FileDialog
@sp 2
@example
@group
r = gDefExt(fd, ext);

object  fd;     /*  file dialog     */
char    *ext;   /*  file extension  */
object  r;      /*  file dialog     */
@end group
@end example
This method is used to set the default file extension which will be
displayed for the user.  It should not contain a period.
@example
@group
@exdent Example:

object  fd;

gDefExt(fd, "txt");
@end group
@end example
@sp 1
See also:  @code{SetFile, AppendFilter}
@end deffn






@deffn {Dispose} Dispose::FileDialog
@sp 2
@example
@group
r = gDispose(fd);

object  fd;   /*  file dialog  */
object  r;    /*  NULL         */
@end group
@end example
This method is used to dispose of a file dialog object.  This method must
be called to dispose of the dialog when it is no longer needed.

The value returned is always @code{NULL} and may be used to null out
the variable which contained the object being disposed in order to
avoid future accidental use.
@example
@group
@exdent Example:

object  fd;

fd = gDispose(fd);
@end group
@end example
@c @sp 1
@c See also:  @code{}
@end deffn











@deffn {GetFile} GetFile::FileDialog
@sp 2
@example
@group
fname = gGetFile(fd);

object  fd;     /*  file dialog  */
char    *fname; /*  file name    */
@end group
@end example
This method is used to obtain the file name the user selected once the
file dialog has been completed.  This value will include the full path,
file name and extension.  If multiple selections are allowed, the path
will be returned, then a space delimited list of selected file names.
@example
@group
@exdent Example:

object  fd;
char    *files;

files = gGetFile(fd);
@end group
@end example
@sp 1
See also:  @code{GetOpenFile, GetSaveFile, SetFile, SetFlags}
@end deffn















@deffn {InitialDir} InitialDir::FileDialog
@sp 2
@example
@group
r = gInitialDir(fd, pth);

object  fd;     /*  file dialog  */
char    *pth;   /*  initial path */
object  r;      /*  file dialog  */
@end group
@end example
This method is used to set the initial path associated with the file dialog.
If no initial path is set, the system will use the current directory of
the application.
@example
@group
@exdent Example:

object  fd;

gInitialDir(fd, "c:\\mypath");
@end group
@end example
@sp 1
See also:  @code{SetFile, AppendFilter, DefExt}
@end deffn







@deffn {New} New::FileDialog
@sp 2
@example
@group
fd = vNew(FileDialog, pwind);

object  pwind;  /*  parent window  */
object  fd;     /*  file dialog    */
@end group
@end example
This method is used to create a new file dialog object.  It will not be
displayed until the appropriate method is called.  The object returned
must be disposed when it no longer needed.

The @code{pwind} parameter refers to either the application's main window
or any child window object.
@example
@group
@exdent Example:

object  fd, pwind;

fd = vNew(FileDialog, pwind);
@end group
@end example
@sp 1
See also:  @code{Dispose, GetOpenFile, GetSaveFile}
@end deffn





@deffn {GetOpenFile} GetOpenFile::FileDialog
@sp 2
@example
@group
r = gGetOpenFile(fd);

object  fd;     /*  file dialog    */
int     r;      /*  return status  */
@end group
@end example
This method is used to actually display the dialog and obtain user input.
The type of file dialog presented will be one in which the user must
select a pre-existing file, presumably to open.

The return value is non-zero if the user makes a valid selection and
zero if the user did not make a selection or if an error occurred.
@example
@group
@exdent Example:

object  fd;
int     r;

r = gGetOpenFile(fd);
@end group
@end example
@sp 1
See also:  @code{GetSaveFile}
@end deffn






@deffn {GetSaveFile} GetSaveFile::FileDialog
@sp 2
@example
@group
r = gGetSaveFile(fd);

object  fd;     /*  file dialog    */
int     r;      /*  return status  */
@end group
@end example
This method is used to actually display the dialog and obtain user input.
The type of file dialog presented will be one in which the user may
select a pre-existing file or type in a new name, presumably to save some
data to.

The return value is non-zero if the user makes a valid selection and
zero if the user did not make a selection or if an error occurred.
@example
@group
@exdent Example:

object  fd;
int     r;

r = gGetSaveFile(fd);
@end group
@end example
@sp 1
See also:  @code{GetOpenFile}
@end deffn








@deffn {SetFile} SetFile::FileDialog
@sp 2
@example
@group
r = gSetFile(fd, fname);

object  fd;     /*  file dialog  */
char    *fname; /*  file name    */
object  r;      /*  file dialog  */
@end group
@end example
This method is used to set the default file name displayed when the
file dialog is displayed.
@example
@group
@exdent Example:

object  fd;

gSetFile(fd, "somefile.txt");
@end group
@end example
@sp 1
See also:  @code{InitialDir, GetFile, SetFlags, DefExt}
@end deffn







@deffn {SetFlags} SetFlags::FileDialog
@sp 2
@example
@group
r = gSetFlags(fd, flgs);

object  fd;     /*  file dialog  */
DWORD   flgs;   /*  flags        */
object  r;      /*  file dialog  */
@end group
@end example
This method is used to set the option flags associated with a file dialog.
The available flags are fully documented in the Windows documentation
under the @code{OPENFILENAME} structure and all begin with @code{OFN_}.
@example
@group
@exdent Example:

object  fd;

gSetFlags(fd, OFN_ALLOWMULTISELECT);
@end group
@end example
@sp 1
See also:  @code{SetFile, DefExt}
@end deffn















@deffn {SetTitle} SetTitle::FileDialog
@sp 2
@example
@group
r = gSetTitle(fd, ttl);

object  fd;     /*  file dialog  */
char    *ttl;   /*  dialog title */
object  r;      /*  file dialog  */
@end group
@end example
This method is used to set the text which appears at the top of the
file dialog.  If no value is set the system will use appropriate defaults.
@example
@group
@exdent Example:

object  fd;

gSetTitle(fd, "Select File");
@end group
@end example
@sp 1
See also:  @code{SetFile, AppendFilter, DefExt}
@end deffn





@subsection Printer Selection and Configuration Dialog
The @code{PrintDialog} class is used to present the user with a standard
and convenient dialog for printer selection and configuration.

This class is seldom needed because the @code{Printer} class has
the @code{QueryPrinter::Printer} method which automatically calls
this class for printer information.  However, this class may be used
independently in order to have better control of user options.


The @code{pd} parameter used by all methods in this class refer to
the print dialog object returned by the @code{New} method.














@deffn {Copies} Copies::PrintDialog
@sp 2
@example
@group
cp = gCopies(fd);

object  fd;     /*  file dialog  */
int     cp;     /*  copies       */
@end group
@end example
This method is used to obtain the number of copies selected by the user
subsequent to calling @code{Perform}.
@example
@group
@exdent Example:

object  fd;
int     cp;

cp = gCopies(fd);
@end group
@end example
@sp 1
See also:  @code{Perform, SetFlags, GetPageRange}
@end deffn









@deffn {Dispose} Dispose::PrintDialog
@sp 2
@example
@group
r = gDispose(pd);

object  pd;     /*  print dialog   */
object  r;      /*  NULL           */
@end group
@end example
This method is used to dispose of a print dialog object.  This must
be done when it is no longer needed.

The value returned is always @code{NULL} and may be used to null out
the variable which contained the object being disposed in order to
avoid future accidental use.
@example
@group
@exdent Example:

object  pd;

pd = gDispose(pd);
@end group
@end example
@sp 1
See also:  @code{New, Perform}
@end deffn











@deffn {GetPageRange} GetPageRange::PrintDialog
@sp 2
@example
@group
r = gGetPageRange(fd, start, end);

object  fd;     /*  file dialog           */
int     *start; /*  starting page number  */
int     *end;   /*  ending page number    */
object  r;      /*  file dialog           */
@end group
@end example
This method is used to obtain the page range selected by the user
subsequent to calling @code{Perform}.
@example
@group
@exdent Example:

object  fd;
int     start, end;

gGetPageRange(fd, &start, &end);
@end group
@end example
@sp 1
See also:  @code{Perform, SetFlags, Copies}
@end deffn











@deffn {Handle} Handle::PrintDialog
@sp 2
@example
@group
hdc = gHandle(pd);

object  pd;     /*  print dialog   */
HANDLE  hdc;    /*  handle to dc   */
@end group
@end example
This method is used to obtain the device context handle associated with
the print dialog.  It will only be valid after performing the
dialog (via @code{Perform}).  If the handle is obtained via this method,
the @code{Dispose} method will not release the handle as it would
normally do.  This is done so the device context can be used for the
printer to be opened.  If the handle as obtained it is the application's
responsibility to release the handle when it is no longer needed.  This
may be done via the @code{DeleteDC} Windows function.

This method is mainly used internally to WDS and is only made available
for convenience.
@example
@group
@exdent Example:

object  pd;
HDC     hdc;

hdc = gHandle(pd);
@end group
@end example
@sp 1
See also:  @code{Perform}
@end deffn









@deffn {New} New::PrintDialog
@sp 2
@example
@group
pd = vNew(PrintDialog, pwind);

object  pwind;  /*  parent window  */
object  pd;     /*  print dialog   */
@end group
@end example
This method is used to create a new print dialog object.  It will not be
displayed until the appropriate method is called.  The object returned
must be disposed when it no longer needed.

The @code{pwind} parameter refers to either the application's main window
or any child window object.
@example
@group
@exdent Example:

object  pd, pwind;

pd = vNew(PrintDialog, pwind);
@end group
@end example
@sp 1
See also:  @code{Perform, Dispose}
@end deffn












@deffn {Perform} Perform::PrintDialog
@sp 2
@example
@group
r = gPerform(pd);

object  pd;     /*  print dialog   */
int     r;      /*  return status  */
@end group
@end example
This method is used to actually present the user with the print dialog.
The return value is non-zero if the user appropriately selects a printer.
Zero will be returned if the user cancels the dialog or an error occurs.
@example
@group
@exdent Example:

object  pd;
int     r;

r = gPerform(pd);
@end group
@end example
@sp 1
See also:  @code{New, Dispose}
@end deffn







@deffn {SetFlags} SetFlags::PrintDialog
@sp 2
@example
@group
r = gSetFlags(pd, flgs);

object  pd;     /*  print dialog  */
DWORD   flgs;   /*  flags         */
object  r;      /*  print dialog  */
@end group
@end example
This method is used to set the option flags associated with a print dialog.
The available flags are fully documented in the Windows documentation
under the @code{PRINTDLG} structure and all begin with @code{PD_}.
@example
@group
@exdent Example:

object  pd;

gSetFlags(pd, PD_ALLPAGES);
@end group
@end example
@c @sp 1
@c See also:  @code{}
@end deffn














@subsection Color Selection Dialog
The @code{ColorDialog} class is used to present the user with a standard
and convenient dialog for color selection and configuration.

The @code{cd} parameter used by all methods in this class refer to
the color dialog object returned by the @code{New} method.









@deffn {Dispose} Dispose::ColorDialog
@sp 2
@example
@group
r = gDispose(cd);

object  cd;     /*  color dialog   */
object  r;      /*  NULL           */
@end group
@end example
This method is used to dispose of a color dialog object.  This must
be done when it is no longer needed.

The value returned is always @code{NULL} and may be used to null out
the variable which contained the object being disposed in order to
avoid future accidental use.
@example
@group
@exdent Example:

object  cd;

cd = gDispose(cd);
@end group
@end example
@sp 1
See also:  @code{New, Perform}
@end deffn









@deffn {GetColor} GetColor::ColorDialog
@sp 2
@example
@group
clr = gGetColor(cd);

object   cd;    /*  color dialog    */
COLORREF clr;   /*  selected color  */
@end group
@end example
This method is used to obtain the users color selection subsequent to
executing @code{Perform}.
@example
@group
@exdent Example:

object  cd;
COLORREF clr;

clr = gGetColor(cd);
@end group
@end example
@sp 1
See also:  @code{Perform, SetColor, GetCustomColorPtr}
@end deffn













@deffn {GetCustomColorPtr} GetCustomColorPtr::ColorDialog
@sp 2
@example
@group
cv = gGetCustomColorPtr(cd);

object   cd;    /*  color dialog    */
COLORREF *cv;   /*  color vector    */
@end group
@end example
This method is used to obtain a pointer to 16 @code{COLORREF}s
selected as custom colors by the user.  This vector should be
inspected subsequent to calling @code{Perform} and prior to
calling @code{Dispose}.  One @code{Dispose} is called, the returned
pointer will no longer be valid.
@example
@group
@exdent Example:

object  cd;
COLORREF *cv;

cv = gGetCustomColorPtr(cd);
@end group
@end example
@sp 1
See also:  @code{Perform, GetColor}
@end deffn




















@deffn {New} New::ColorDialog
@sp 2
@example
@group
cd = vNew(ColorDialog, pwind);

object  pwind;  /*  parent window  */
object  cd;     /*  color dialog   */
@end group
@end example
This method is used to create a new color dialog object.  It will not be
displayed until the appropriate method is called.  The object returned
must be disposed when it no longer needed.

The @code{pwind} parameter refers to either the application's main window
or any child window object.
@example
@group
@exdent Example:

object  cd, pwind;

cd = vNew(ColorDialog, pwind);
@end group
@end example
@sp 1
See also:  @code{Perform, Dispose}
@end deffn









@deffn {Perform} Perform::ColorDialog
@sp 2
@example
@group
r = gPerform(cd);

object  cd;     /*  color dialog   */
int     r;      /*  return status  */
@end group
@end example
This method is used to actually present the user with the color
selection dialog.  The return value is non-zero if the user appropriately
selects a color.  Zero will be returned if the user canceled the dialog
or an error occurred.
@example
@group
@exdent Example:

object  cd;
int     r;

r = gPerform(cd);
@end group
@end example
@sp 1
See also:  @code{New, Dispose}
@end deffn









@deffn {SetColor} SetColor::ColorDialog
@sp 2
@example
@group
r = gSetColor(cd, clr);

object   cd;    /*  color dialog    */
COLORREF clr;   /*  selected color  */
object   r;     /*  color dialog    */
@end group
@end example
This method is used to set the initial color associated with the
color selection dialog.  The Windows @code{RGB} macro may be used
to obtain a valid @code{COLORREF}.
@example
@group
@exdent Example:

object  cd;
COLORREF clr;

gSetColor(cd, RGB(255, 0 ,0));
@end group
@end example
@sp 1
See also:  @code{Perform, GetColor}
@end deffn














@deffn {SetFlags} SetFlags::ColorDialog
@sp 2
@example
@group
r = gSetFlags(cd, flgs);

object  cd;     /*  color dialog  */
DWORD   flgs;   /*  flags         */
object  r;      /*  color dialog  */
@end group
@end example
This method is used to set the option flags associated with a color dialog.
The available flags are fully documented in the Windows documentation
under the @code{CHOOSECOLOR} structure and all begin with @code{CC_}.
@example
@group
@exdent Example:

object  cd;

gSetFlags(cd, CC_FULLOPEN);
@end group
@end example
@c @sp 1
@c See also:  @code{}
@end deffn











@subsection Font Selection Dialog
The @code{FontDialog} class is used to present the user with a standard
and convenient dialog for font selection.

The @code{fd} parameter used by all methods in this class refer to
the font dialog object returned by the @code{New} method.








@deffn {Dispose} Dispose::FontDialog
@sp 2
@example
@group
r = gDispose(fd);

object  fd;     /*  font dialog   */
object  r;      /*  NULL          */
@end group
@end example
This method is used to dispose of a font dialog object.  This must
be done when it is no longer needed.

The value returned is always @code{NULL} and may be used to null out
the variable which contained the object being disposed in order to
avoid future accidental use.
@example
@group
@exdent Example:

object  fd;

fd = gDispose(fd);
@end group
@end example
@sp 1
See also:  @code{New, Perform}
@end deffn











@deffn {Font} Font::FontDialog
@sp 2
@example
@group
fnt = gFont(fd);

object  fd;     /*  font dialog  */
object  font;   /*  font object  */
@end group
@end example
This method is used to create a font object representing the font the
user selected.  This font object will be an instance of the
@code{ExternalFont} class and must be disposed (either directly or
indirectly) when it is no longer needed.  This method should only
be called subsequent to @code{Perform}.  @code{NULL} will be returned
if no font was selected or the font object could not be created.
@example
@group
@exdent Example:

object  fd, fnt;

fnt = gFont(fd);
@end group
@end example
@sp 1
See also:  @code{Perform}
@end deffn











@deffn {GetColor} GetColor::FontDialog
@sp 2
@example
@group
clr = gGetColor(fd);

object   fd;    /*  font dialog  */
COLORREF clr;   /*  color        */
@end group
@end example
This method is used to get the user selected color once @code{Perform}
has been called.  This will only work if the @code{CF_EFFECTS} option
is set.
@example
@group
@exdent Example:

object  fd;
COLORREF clr;

clr = gGetColor(fd);
@end group
@end example
@sp 1
See also:  @code{SetFlags, SetColor, Perform}
@end deffn














@deffn {New} New::FontDialog
@sp 2
@example
@group
fd = vNew(FontDialog, pwind);

object  pwind;  /*  parent window  */
object  fd;     /*  font dialog    */
@end group
@end example
This method is used to create a new font dialog object.  It will not be
displayed until the appropriate method is called.  The object returned
must be disposed when it no longer needed.

The @code{pwind} parameter refers to either the application's main window
or any child window object.
@example
@group
@exdent Example:

object  fd, pwind;

fd = vNew(FontDialog, pwind);
@end group
@end example
@sp 1
See also:  @code{Perform, Dispose}
@end deffn









@deffn {Perform} Perform::FontDialog
@sp 2
@example
@group
r = gPerform(fd);

object  fd;     /*  font dialog    */
int     r;      /*  return status  */
@end group
@end example
This method is used to actually present the user with the font
selection dialog.  The return value is non-zero if the user appropriately
selects a font.  Zero will be returned if the user canceled the dialog
or an error occurred.
@example
@group
@exdent Example:

object  fd;
int     r;

r = gPerform(fd);
@end group
@end example
@sp 1
See also:  @code{New, Dispose}
@end deffn









@deffn {SetColor} SetColor::FontDialog
@sp 2
@example
@group
clr = gSetColor(fd);

object   fd;    /*  font dialog  */
COLORREF clr;   /*  color        */
@end group
@end example
This method is used to set the initial color of the user color selection
portion of the font dialog.  This will only work if the
@code{CF_EFFECTS} option is set.  The Windows supplied @code{RGB} macro
may be used to create the @code{COLORREF}.
@example
@group
@exdent Example:

object  fd;

gSetColor(fd, RGB(255, 0, 0));
@end group
@end example
@sp 1
See also:  @code{SetFlags, GetColor, Perform}
@end deffn











@deffn {SetFlags} SetFlags::FontDialog
@sp 2
@example
@group
r = gSetFlags(fd, flgs);

object  fd;     /*  font dialog  */
DWORD   flgs;   /*  flags        */
object  r;      /*  font dialog  */
@end group
@end example
This method is used to set the option flags associated with a font dialog.
The available flags are fully documented in the Windows documentation
under the @code{CHOOSEFONT} structure and all begin with @code{CF_}.
@example
@group
@exdent Example:

object  fd;

gSetFlags(fd, CF_ANSIONLY);
@end group
@end example
@c @sp 1
@c See also:  @code{}
@end deffn






@section Dynamic Link Libraries
The @code{DynamicLibrary} class is used to support the dynamic link
library (DLL) facility of Windows.  This class provides a convenient
mechanism to load, free and use DLLs.

Throughout this section @code{dl} will refer to the object which was
returned by the @code{LoadLibrary} method and represents a DLL.








@deffn {DeepDispose} DeepDispose::DynamicLibrary
@sp 2
This method performs the same function as @code{Dispose}.  See that
method for details.
@end deffn







@deffn {Dispose} Dispose::DynamicLibrary
@sp 2
@example
@group
r = gDispose(dl);

object  dl;   /*  a DLL object  */
object  r;    /*  NULL          */
@end group
@end example
This method is used to release and dispose of a DLL object when it is no
longer needed.  All DLL objects are automatically released by WDS when
the application terminates.

The value returned is always @code{NULL} and may be used to null out
the variable which contained the object being disposed in order to
avoid future accidental use.
@example
@group
@exdent Example:

object  dl;

dl = gDispose(dl);
@end group
@end example
@sp 1
See also:  @code{FreeAll, DeepDispose}
@end deffn









@deffn {FindStr} FindStr::DynamicLibrary
@sp 2
@example
@group
dl = gFindStr(DynamicLibrary, fname);

char    *fname; /*  DLL file name  */
object  dl;     /*  DLL object     */
@end group
@end example
This class method is used to obtain the DLL object (previously loaded via
@iftex
@break 
@end iftex
@code{LoadLibrary}) via its associated file name.  This won't work
after the DLL object is disposed.
@example
@group
@exdent Example:

object  dl;

dl = gFindStr(DynamicLibrary, "mydll.dll");
@end group
@end example
@sp 1
See also:  @code{LoadLibrary}
@end deffn










@deffn {FreeAll} FreeAll::DynamicLibrary
@sp 2
@example
@group
gFreeAll(DynamicLibrary);

@end group
@end example
This class method is used to release and dispose of all DLL objects.
It is automatically called by WDS when the application terminates.
@example
@group
@exdent Example:

gFreeAll(DynamicLibrary);
@end group
@end example
@sp 1
See also:  @code{Dispose, DeepDispose}
@end deffn











@deffn {GetProcAddress} GetProcAddress::DynamicLibrary
@sp 2
@example
@group
fun = gGetProcAddress(dl, fname);

object  dl;     /*  DLL object           */
char    *fname; /*  function name        */
FARPROC fun;    /*  pointer to function  */
@end group
@end example
This method is used to obtain a pointer to a function within a
DLL.  The DLL function may then be evoked with the pointer.
@example
@group
@exdent Example:

FARPROC fun;
object  dl;

fun = gGetProcAddress(dl, "myfunc");
@end group
@end example
@sp 1
See also:  @code{FindStr, LoadLibrary}
@end deffn










@deffn {LoadLibrary} LoadLibrary::DynamicLibrary
@sp 2
@example
@group
dl = gLoadLibrary(DynamicLibrary, file);

char    *file;  /*  DLL file    */
object  dl;     /*  DLL object  */
@end group
@end example
This method is used to open a DLL file and create a WDS object which
represents the loaded library.  If the named file does not contain a
path, the normal Windows search path will be searched.  If the
DLL cannot be found or loaded, @code{NULL} will be returned.
@example
@group
@exdent Example:

object  dl;

dl = gLoadLibrary(DynamicLibrary, "mydll.dll");
@end group
@end example
@sp 1
See also:  @code{GetProcAddress, Dispose, FindStr}
@end deffn








